\section{Limitations and Open Engineering Problems}
\label{sec:limitations}

Argus gives an agent real judgment, real tool access, and real persistence. It
does not guarantee good science, and a report that documented reliability
mechanisms without documenting their boundaries would itself be misleading. The
limitations below are material, and several are open engineering problems the
roadmap (Section~\ref{sec:roadmap}) targets directly.

\subsection{Model and objective ceilings}

\begin{description}[leftmargin=1.4em, style=nextline]
  \item[Provider dependence.] Argus drives an external agent CLI (Codex, Claude
    Code, or Copilot CLI) and inherits that provider's capability, availability,
    latency, and pricing. An outage, a regression, or a rate limit is felt
    directly, and the quality ceiling of any run is the underlying model's.
  \item[Objective quality is the operator's.] The system optimizes the objective
    it is given. A vague or trivially gameable objective yields vague or gamed
    work; Argus enforces integrity of \emph{measurement}, not wisdom of
    \emph{goal}.
  \item[No guarantee of novelty, correctness, or SOTA.] Nothing in the runtime
    guarantees that an idea is novel, that a proof or result is correct, or that a
    number is state of the art. Those are properties of the agent's judgment and
    the problem's difficulty, and the system deliberately does not fake them.
\end{description}

\subsection{Sole-reviewer completion authority}

The Reviewer is the single completion authority, and it is a model. A wrong
Reviewer can accept an incomplete result or reject a finished one. Argus reduces
this risk by giving the Reviewer the shared work-tree, the execution-log audit, and
a structured checklist (Section~\ref{sec:evidence}), and by routing a degraded
reviewer backend to an explicit infrastructure-failure status rather than a
\code{continue}. But it cannot eliminate model error, and there is no independent
oracle behind the Reviewer. The historical blind-gate incident
(Section~\ref{sec:reliability}), in which a stale schema path ran the completion
gate blind for roughly ninety minutes, is the concrete failure mode of concentrating
authority in one component; the fix hardened the transport, not the underlying
concentration of authority, which remains an open problem.

\subsection{Evidence scope and integrity coverage}

\begin{description}[leftmargin=1.4em, style=nextline]
  \item[Four of six results are snapshots.] Only two of the six public results
    (Section~\ref{sec:results}) carry committed digest records referencing
    artifacts in named external projects; the artifact bytes are not in this
    repository. The other four are public website snapshots linked to the evidence
    bundle. They are real, scoped claims under their protocols, but they are not
    yet independently reproducible from a named commit, and the report does not
    present them as if they were.
  \item[Integrity is per-benchmark and ongoing.] The anti-cheat construction is
    task-agnostic, but a determined agent explores the exploit surface of each
    specific benchmark. A new benchmark can present a new exploit --- a leaked test
    set, an unrandomized input, a side channel --- that the current checks were not
    written for. Measurement integrity is an ongoing engineering surface, not a
    solved problem.
  \item[Paper inventory is not acceptances.] The 41-paper portfolio
    (Section~\ref{sec:portfolio}) is a de-duplicated inventory compared to
    human-authored literature; it is not a count of accepted papers, and no
    acceptance is claimed.
\end{description}

\subsection{Operational limitations}

\begin{description}[leftmargin=1.4em, style=nextline]
  \item[Cost.] Long-horizon autonomy consumes provider tokens continuously.
    Budgets and reservations bound spend, but 7$\times$24 operation is not free.
  \item[Secret handling is a scanner, not a proof.] The credential guard scrubs
    artifacts before they are recorded, but it is a task-agnostic scanner: a scan
    gap or an oversized artifact blocks certification rather than guaranteeing zero
    leakage, and operators remain responsible for how secrets enter the environment.
  \item[Back-compatibility surface.] Legacy shims (an older global-directory
    variable and queue helpers) are kept for migration; they are technical debt to
    retire and add surface an audit must account for.
  \item[GPU lease is invoked, not automatic.] The GPU lease
    (Section~\ref{sec:execution}) is robust but is an operator/agent-invoked tool,
    not an automatic runtime service; coordinating scarce hardware still depends on
    the agent choosing to use it.
\end{description}

\subsection{Responsible operation}

Given these limitations, Argus is meant to be operated and audited, not fired and
forgotten. Responsible use means giving it well-posed objectives and honest
benchmarks; setting budgets deliberately; reading the Reviewer's reasoning rather
than a headline; treating every number according to its evidence tier
(Table~\ref{tab:tiers}); and publishing only what the stated evidence supports. The
system makes this practical --- persisted artifacts, a round-by-round audit trail, a
single inspectable completion authority --- but the discipline of honest reporting
ultimately rests with the operator.

\begin{callout}[Maturity]
Argus is under active development at Technical Report 0.3. The architecture and
mechanisms described here are implemented and grounded in the current source tree;
the public results (Section~\ref{sec:results}) are scoped claims under their stated
protocols and evidence tiers, not universal-capability guarantees. Treat any
performance number according to its tier --- an artifact-digest record, a
public website snapshot linked to the bundle, or a labeled external baseline ---
never as a general guarantee of capability.
\end{callout}
